% =============================================================================
% APÊNDICE — Aparato de Leitura
% Estudo crítico pós-ficção. Entra no \backmatter, DEPOIS da Nota Histórica e
% ANTES da contracapa. Projeto de imersão: a ficção termina primeiro; um
% LIMIAR em sépia (última voz do Arquivista) marca a saída; o aparato roda em
% PAPEL BRANCO, registro sóbrio — contraste que protege (e valoriza) o peso
% sensorial do livro, em vez de diluí-lo.
% =============================================================================

% Acento local (evita depender de cores definidas só no dossiê externo).
\definecolor{aparatoacc}{HTML}{7A5C34}

% -----------------------------------------------------------------------------
% LIMIAR DE SAÍDA — ainda no sépia, ainda em-diegese: a passagem.
% -----------------------------------------------------------------------------
\cleardoublepage
\thispagestyle{empty}
\vspace*{\fill}
\begin{center}
{\scshape\large O arquivo termina aqui.}\\[1.6em]

\begin{minipage}{0.82\textwidth}\itshape\centering
O que segue não pertence ao arquivo.\\[0.6em]
São notas de quem ficou depois que a luz se apagou --- e quis entender
\textup{como} o medo foi construído: a máquina por trás do cheiro.\\[0.6em]
Se você prefere guardar o susto inteiro, feche o livro agora.\\[0.3em]
Ainda dá tempo.\\[0.3em]
Quem virar a página estará saindo do escuro por conta própria.
\end{minipage}\\[2.2em]

{\scshape --- O Arquivista}
\end{center}
\vspace*{\fill}
\clearpage

% -----------------------------------------------------------------------------
% VIRADA DE REGISTRO: papel branco, tinta neutra, cabeçalho sóbrio.
% A partir daqui o leitor saiu do arquivo — o livro assume o tom de estudo.
% -----------------------------------------------------------------------------
\cleardoublepage
\pagecolor{white}
\color{black}
\pagestyle{plain}

\chapter*{Apêndice --- Aparato de Leitura}
\addcontentsline{toc}{chapter}{Apêndice --- Aparato de Leitura}
\markboth{Aparato de Leitura}{}

\noindent\textit{Um estudo do avesso da obra.} As páginas anteriores quiseram
fazer você \emph{sentir}; estas querem mostrar \emph{como}. É material de
consulta para leitores curiosos, estudantes de literatura e professores ---
e pode ser lido, como o resto do livro, na ordem que quiser.

\vspace{0.4em}
\noindent{\small\color{aparatoacc}\rule{\linewidth}{0.6pt}}

\vspace{0.6em}
\noindent\textbf{Três honestidades.} \emph{(1)} Isto não é biografia do autor:
o que se lê aqui é a obra e o seu método, não a vida de quem a escreveu.
\emph{(2)} As \emph{influências} são de dois tipos --- \textbf{fonte citada}
(que a obra nomeia) e \textbf{afinidade} (parentesco que o crítico
reconhece) --- e o texto sempre diz qual é qual. \emph{(3)} Não se declara a
obra ``clássico'' nem ``obra-prima'': isso dependeria de recepção e validação
externas que ela ainda não teve. Afirma-se o que se pode ler no texto --- nada
além.

% =============================================================================
\section*{A. Estilo e prosa}
\addcontentsline{toc}{section}{A. Estilo e prosa}
\markboth{Aparato --- Estilo e prosa}{}

Se a estrutura é o esqueleto do livro, o \textbf{estilo} é a pele que faz o
horror ser sentido. A prosa é feita de registros que se atritam de propósito.

\subsection*{Um estilo de três registros}
\begin{itemize}
\item \textbf{Documental-frio} (fragmentos DOC): sintaxe de laudo, impessoal;
o horror vaza por \emph{anomalia} --- rasura, lacuna, troca de termo.
\item \textbf{Confessional-lírico} (MEM/CONT): primeira e segunda pessoa,
imagens sensoriais, ritmo de intimidade e de cerco.
\item \textbf{Forense-profissional} (LUC): a perita que raciocina em voz alta,
entre o método científico e a contaminação.
\end{itemize}
\noindent A \emph{fricção} entre eles --- um laudo seco descrevendo o
inominável --- é o motor do estilo: a \emph{dissonância clínica\slash narrativa}.

\subsection*{Sintaxe: parataxe, fragmento e anáfora}
Predomina a frase curta, declarativa, muitas vezes reduzida a \emph{fragmento}
(``Mas não. Você abriu.'') --- respiração ofegante, de perseguição. A
\textbf{anáfora} organiza blocos inteiros (``Podia ter\ldots\ Podia ter\ldots'';
``Recomenda-se que\ldots''), criando um ritmo de sentença inescapável. O
\textbf{assíndeto} acelera; o travessão e as reticências pontuam o silêncio.

\subsection*{Léxico sensorial e o corpo}
O vocabulário é corporal e olfativo: o ``cheiro adocicado'', a carne, o calor
que sobe ``entre os órgãos'', a brasa no olho. Em vez de \emph{explicar} o
medo, o texto o faz salivar, arder, pesar --- é a \emph{ancoragem sensorial}.

\subsection*{Repetição, refrão e \emph{leitmotiv}}
Sintagmas retornam como refrões (``a fome''; ``o ciclo recomeça a cada
leitor''), à maneira de um \emph{leitmotiv} musical. Para não virar maneirismo,
usa-se a \textbf{repetição quebrada}: manter o refrão, mas interrompê-lo antes
do esperado ou deslocá-lo para outro sujeito.

\subsection*{Diálogo, aforismo, tempo e materialidade}
Nas sessões (LUC-08), o diálogo torna-se \textbf{esticomitia} --- réplicas
curtíssimas, um duelo verbal em que o não-dito faz o medo. A voz aforística
condensa a tese em \textbf{antítese} (``Ela não é um monstro. Ela é uma
consequência.''), e o Arquivista fala por imperativos. O \textbf{presente
intrusivo} dos CONT (``agora você está lendo esta linha'') produz cerco,
contra o pretérito memorialístico das MEM. E a \textbf{materialidade
tipográfica} --- cor por tipo de fragmento, marginália manuscrita, capitular,
sépia --- é estilo: o suporte significa junto com a frase.

% =============================================================================
\section*{B. Estudo para acadêmicos de literatura}
\addcontentsline{toc}{section}{B. Estudo para acadêmicos de literatura}
\markboth{Aparato --- Estudo acadêmico}{}

\noindent\textit{Enquadramento honesto:} as molduras abaixo são \textbf{lentes
críticas} para \emph{ler} a obra --- não são método declarado pelo autor nem
prova de que ele tenha lido esses teóricos.

\subsection*{Molduras teóricas aplicáveis}
\begin{description}[style=nextline,leftmargin=0pt,itemsep=3pt]
\item[O fantástico como hesitação --- Todorov (1970)] A recusa sustentada de
confirmar a criatura mantém o texto na \emph{hesitação} entre o estranho e o
maravilhoso: o fantástico ``puro''.
\item[Narrador não confiável --- Booth (1961); Chatman] Joaquim, interno
psiquiátrico: o discurso mina a própria voz que narra.
\item[Focalização e metalepse --- Genette] A interpelação em segunda pessoa
(CONT) cruza os níveis narrativos --- o leitor é puxado para dentro da diegese.
\item[Paratexto e transtextualidade --- Genette, \textit{Seuils} (1987)] A
Nota do Arquivista, o ``Como Ler'', a ficha e a marginália são
\textbf{paratexto ficcionalizado}: a moldura vira ficção.
\item[Recepção: leitor implícito e vazios --- Iser (1976)] Os \emph{vazios}
(\emph{Leerstellen}) deliberados: o leitor \emph{completa} o horror.
\item[Literatura ergódica / cibertexto --- Aarseth (1997)] As rotas de leitura
exigem ``esforço não trivial'' para percorrer o texto: hipertexto em papel.
\item[Metaficção historiográfica --- Hutcheon (1988)] A ficção expõe sua
relação com a história documentada (Barbacena) e problematiza o arquivo.
\item[Trauma e latência --- Caruth (1996)] A ``fome'' como trauma não
assimilado que retorna a destempo.
\item[Pós-memória --- Hirsch (2012)] A transmissão do trauma a quem não o
viveu: o leitor herda como Paciente 1.263.
\item[Abjeção --- Kristeva (1980)] O cadáver, a vala, o cheiro, o corpo em
ruína: o horror da fronteira corporal desfeita.
\item[Filosofia do horror --- Carroll (1990)] O ``monstro'' como violação
categorial (real\slash mental\slash textual) sustenta o \emph{art-horror}.
\item[Polifonia e heteroglossia --- Bakhtin] O mosaico de vozes ---
institucional, íntima, forense --- como romance polifônico.
\item[Topos do manuscrito encontrado --- Cervantes; gótico; Eco] O diário
``achado'' no peito do morto: dispositivo de autenticação e de dúvida.
\end{description}

\subsection*{Agenda de pesquisa (recortes possíveis)}
\begin{enumerate}[label=\textbf{\alph*)},leftmargin=1.6em]
\item A ambiguidade como método: Molambudos e o fantástico ``puro'' de Todorov.
\item Paratexto ficcionalizado: a moldura do Arquivista à luz de Genette.
\item Do arquivo ao leitor: pós-memória e a cadeia de hospedeiros.
\item Cibertexto impresso: as rotas de leitura como literatura ergódica.
\item Ética da representação: metaficção historiográfica e o Hospital Colônia.
\item A abjeção sertaneja: corpo, fome e vala à luz de Kristeva.
\end{enumerate}

\subsection*{Bibliografia teórica de apoio}
{\small
AARSETH, E. \textit{Cybertext}. 1997. \textbullet\
BAKHTIN, M. \textit{Problemas da Poética de Dostoiévski}. \textbullet\
BOOTH, W. \textit{A Retórica da Ficção}. 1961. \textbullet\
CARROLL, N. \textit{A Filosofia do Horror}. 1990. \textbullet\
CARUTH, C. \textit{Unclaimed Experience}. 1996. \textbullet\
GENETTE, G. \textit{Paratextos Editoriais}. 1987. \textbullet\
HIRSCH, M. \textit{The Generation of Postmemory}. 2012. \textbullet\
HUTCHEON, L. \textit{Poética do Pós-Modernismo}. 1988. \textbullet\
ISER, W. \textit{O Ato da Leitura}. 1976. \textbullet\
KRISTEVA, J. \textit{Poderes do Horror}. 1980. \textbullet\
TODOROV, T. \textit{Introdução à Literatura Fantástica}. 1970.\par
\vspace{0.3em}\noindent\emph{Base histórica (fonte citada na obra):}
ARBEX, D. \textit{Holocausto Brasileiro}. 2013.}

% =============================================================================
\section*{C. Dossiê para estudantes de literatura contemporânea}
\addcontentsline{toc}{section}{C. Dossiê para estudantes de literatura contemporânea}
\markboth{Aparato --- Dossiê para estudantes}{}

\subsection*{A obra em seu tempo}
\textit{Molambudos} dialoga com correntes bem reconhecíveis do presente
literário --- sem que isso a torne ``pioneira'' (o que seria exagero):
\begin{itemize}
\item \textbf{Virada da memória e do trauma:} a literatura e o jornalismo
voltaram-se para os mortos sem registro; a repercussão de \textit{Holocausto
Brasileiro} (2013) é parte desse movimento de reparação pela memória.
\item \textbf{Ficção documental (não autoficção):} o livro escolhe o
\emph{documento} --- prontuário, laudo, termo --- para autenticar o horror
pela frieza do arquivo.
\item \textbf{Sensibilidade digital e hipertexto:} rotas de leitura e ``texto
que contamina'' ecoam a cultura de wikis, arquivos e \emph{creepypasta}\slash
SCP, aqui materializada em papel.
\item \textbf{Horror como literatura legítima:} o medo como via para falar de
luto, classe e trauma (o chamado \emph{elevated horror}).
\item \textbf{O sertão reativado:} a grande tradição da seca relida pela lente
do horror e da memória institucional.
\end{itemize}

\subsection*{Catálogo de técnicas literárias (para reconhecer no texto)}
Uma lista de consulta: a técnica, o que é, e onde procurá-la no livro.
\begin{description}[style=nextline,leftmargin=0pt,itemsep=3pt]
\item[Narrativa fragmentária] O sentido nasce da montagem de peças
independentes. \emph{Onde:} os 67 fragmentos em 4 tipos.
\item[Ficção documental (\emph{found document})] Ficção disfarçada de documento
real. \emph{Onde:} os fragmentos DOC.
\item[Narrador não confiável] Voz cuja credibilidade o texto põe em dúvida.
\emph{Onde:} Joaquim.
\item[Metaficção e segunda pessoa] A ficção comenta o próprio ato de ler e
interpela ``você''. \emph{Onde:} os CONT.
\item[Metalepse] Quebra da fronteira entre narração e leitor. \emph{Onde:} o
leitor vira Paciente 1.263.
\item[Dissonância clínica\slash narrativa] Linguagem burocrática fria sobre
conteúdo atroz. \emph{Onde:} os laudos.
\item[Ancoragem sensorial] Trocar explicação por resposta corporal.
\emph{Onde:} o cheiro, o calor, a fome.
\item[Falsa segurança] Fingir que o perigo cessou --- e quebrar a calma no
mesmo parágrafo. \emph{Onde:} aberturas de vários CONT.
\item[Repetição quebrada] Manter um refrão e interrompê-lo antes do esperado.
\emph{Onde:} ``a fome'', ``o ciclo''.
\item[Lacuna sugestiva] Deixar o leitor completar o horror por ausência.
\emph{Onde:} frases cortadas, documentos com falhas.
\item[Esticomitia] Diálogo de réplicas curtíssimas. \emph{Onde:} LUC-08.
\item[Paratexto ficcionalizado] Nota, ficha e marginália como parte da ficção.
\emph{Onde:} a moldura do Arquivista.
\item[Ambiguidade (não indecisão)] O texto \emph{omite} e deixa você decidir,
em vez de afirmar as duas coisas. \emph{Onde:} a natureza da criatura.
\end{description}

\subsection*{Cinco perguntas para discussão}
\begin{enumerate}[leftmargin=1.6em]
\item Por que o livro começa pela fome \emph{real} antes de qualquer
sobrenatural?
\item Onde o \emph{documento} começa a ``falhar'' --- e o que isso provoca?
\item Em que ponto você deixou de confiar no narrador?
\item O livro chamar você de ``Paciente 1.263'' é manipulação ou arte do
suspense?
\item A criatura é entidade, trauma ou texto que contamina? Defenda com
citações --- e note se o livro \emph{permite} uma resposta única.
\end{enumerate}

\vspace{0.6em}
\noindent{\small\color{aparatoacc}\rule{\linewidth}{0.6pt}}
\vspace{0.4em}

\noindent\small\textit{Nota de responsabilidade.} Esta é obra para maiores de
18 anos, apoiada em tragédia histórica real (os currais do Ceará; o Hospital
Colônia de Barbacena). Em sala de aula, recomenda-se contextualizar o material
histórico antes da leitura e tratar o horror como \emph{via de memória} ---
nunca como choque gratuito. O pior que se pode fazer com os mortos, diz a Nota
Histórica, não é assombrá-los com uma história: é deixá-los sem nenhuma.

\clearpage
